\documentclass[../../main.tex]{subfiles}
\graphicspath{{../../image/}{../../figure/}}
\begin{document}

\section{双曲性与严格双曲性}\label{section:ch07-双曲性}

如前几章一样，为避免非本质的技术困难，理论将在典范形式的齐次守恒律组的框架内展开：
\begin{align}
\partial_t U(x,t)+\partial_x F(U(x,t))=0.
\label{align:ch07-典范守恒律}
\end{align}
$F$ 是从 $\mathbb{R}^n$ 的开凸子集 $\mathscr{O}$ 到 $\mathbb{R}^n$ 的 $C^3$ 映射。

在应用中，方程组\eqref{align:ch07-典范守恒律}常控制平面前解，即 $U=U(\nu\cdot x,t)$，空间方向 $\nu\in\mathbb{S}^{m-1}$，来自多维守恒律组\eqref{align:ch04-守恒律组}
。与此相关，
\begin{align}
F(U)=\sum_{\alpha=1}^m\nu_\alpha G_\alpha(U),\qquad U\in\mathscr{O}.
\label{align:ch07-平面前F}
\end{align}

参考第 7.1 节引入的例子，为把热弹性方程组\eqref{align:ch07-热弹性}写成典范形式，我们必须从 $(u,v,s)$ 切换到新状态变量 $(u,v,E)$，其中 $E=\varepsilon+\frac12 v^2$ 是总能量。类似地，等熵气体流动的方程组\eqref{align:ch07-等熵Euler}以状态变量 $(\rho,m)$ 写成典范形式，其中 $m=\rho v$ 是动量。

由\cref{definition:ch03-沿t方向双曲性}
，方程组\eqref{align:ch07-典范守恒律}是双曲的，如果对每一 $U\in\mathscr{O}$，$n\times n$ Jacobi 矩阵 $\mathrm{D}F(U)$ 具有实特征值 $\lambda_1(U)\leqslant\cdots\leqslant\lambda_n(U)$ 与 $n$ 个线性无关的特征向量 $R_1(U),\cdots,R_n(U)$。为后文使用，我们还引入 $\mathrm{D}F(U)$ 的左（行）特征向量 $L_1(U),\cdots,L_n(U)$，规范化为
\begin{align}
L_i(U)R_j(U)=
\begin{cases}
0,& i\neq j\\
1,& i=j.
\end{cases}
\label{align:ch07-左右向量规范化}
\end{align}
此后，符号 $\lambda_i$、$R_i$ 与 $L_i$ 将专门用于表示这些对象。

显然，多维方程组\eqref{align:ch04-守恒律组}
是双曲的，当且仅当通过\eqref{align:ch07-平面前F}由它产生的一切一维方程组\eqref{align:ch07-典范守恒律}（对任意 $\nu\in\mathbb{S}^{m-1}$）都是双曲的。因而双曲性本质上是一维概念。

对于 Lagrange 坐标下一维等熵弹性的方程组\eqref{align:ch07-p系统}（它将在整章中作为阐释一般概念的载体），我们有
\begin{align}
\lambda_1=-\sigma'(u)^{1/2},\qquad \lambda_2=\sigma'(u)^{1/2},
\label{align:ch07-p系统特征值}
\end{align}
\begin{align}
R_1=\tfrac12\begin{pmatrix}-\sigma'(u)^{-1/2}\\-1\end{pmatrix},\qquad
R_2=\tfrac12\begin{pmatrix}-\sigma'(u)^{-1/2}\\1\end{pmatrix},
\label{align:ch07-p系统右特征}
\end{align}
\begin{align}
L_1=(-\sigma'(u)^{1/2},-1),\qquad L_2=(-\sigma'(u)^{1/2},1).
\label{align:ch07-p系统左特征}
\end{align}

$\mathrm{D}F$ 的特征值 $\lambda_i$（$i=1,\cdots,n$）称为方程组\eqref{align:ch07-典范守恒律}的\emph{$i$-特征速度}。该术语源于下述内容。

\begin{definition}[$i$-特征线]\label{definition:ch07-i特征线}
方程组\eqref{align:ch07-典范守恒律}的一条\emph{$i$-特征线}（$i=1,\cdots,n$），关联于一个古典解 $U$，是一条 $C^1$ 函数 $x=x(t)$，其图像含于 $U$ 的定义域内，且它是常微分方程
\begin{align}
\frac{\mathrm{d}x}{\mathrm{d}t}=\lambda_i(U(x,t)).
\label{align:ch07-特征线ODE}
\end{align}
的积分曲线。
\end{definition}

常微分方程\eqref{align:ch07-特征线ODE}的标准存在唯一性理论蕴含：通过方程组\eqref{align:ch07-典范守恒律}的古典解定义域内任一点 $(\bar{x},\bar{t})$，每一特征线族恰有一条特征线穿过。

特征线是各类波的载体。例如，一般平衡律方程组\eqref{align:ch01-closed-system}
在特化到\eqref{align:ch07-典范守恒律}时的方程\eqref{align:ch01-weak-front-jump}
蕴含弱前沿沿特征线传播。因而，$\mathrm{D}F$ 多重特征值的存在可能因共振而在解的行为中诱发严重的复杂性。于是自然要把那些摆脱此类复杂的方程组单独挑出：

\begin{definition}[严格双曲]\label{definition:ch07-严格双曲}
方程组\eqref{align:ch07-典范守恒律}是\emph{严格双曲的}，如果对任意 $U\in\mathscr{O}$，Jacobi 矩阵 $\mathrm{D}F(U)$ 具有相异的实特征值
\begin{align}
\lambda_1(U)<\cdots<\lambda_n(U).
\label{align:ch07-严格双曲不等式}
\end{align}
\end{definition}

由\eqref{align:ch07-p系统特征值}，Lagrange 坐标下等熵弹性的方程组\eqref{align:ch07-p系统}是严格双曲的。绝热热弹性方程组\eqref{align:ch07-热弹性}同理是严格双曲的，其特征速度为
\begin{align}
\lambda_1=-\sqrt{-p_u(u,s)},\qquad \lambda_2=0,\qquad \lambda_3=\sqrt{-p_u(u,s)}.
\label{align:ch07-热弹性特征速度}
\end{align}

理想气体的方程组\eqref{align:ch07-理想气体}具有特征速度
\begin{align}
\lambda_1=v-(\kappa\gamma)^{1/2}\rho^{\frac{\gamma-1}{2}},\qquad
\lambda_2=v+(\kappa\gamma)^{1/2}\rho^{\frac{\gamma-1}{2}},
\label{align:ch07-理想气体特征速度}
\end{align}
因而它在状态空间中 $\rho>0$ 的部分上是严格双曲的。

进而，由二维等熵流动的 Euler 方程通过\eqref{align:ch07-平面前F}产生的任一一维方程组都是严格双曲的。

鉴于上述例子，读者可能形成这样的印象：严格双曲性是连续介质物理中方程组的常态。然而事实并非如此。例如，平面弹性振荡的方程组\eqref{align:ch07-平面振荡}在那些方向 $\nu$（声学张量\eqref{align:ch03-声张量}
在其中具有多重特征值）上不是严格双曲的。事实上，已经证明：在由母三维方程组\eqref{align:ch04-守恒律组}
通过\eqref{align:ch07-平面前F}产生的规模 $n=\pm2,\pm3,\pm4\pmod 8$ 的一维方程组\eqref{align:ch07-典范守恒律}中，严格双曲性至少在某个空间方向 $\nu\in\mathbb{S}^2$ 上必然失效。特别地，由二维非等熵流动的 Euler 方程产生的一维方程组（$n=4$），或由三维等熵或非等熵流动产生的一维方程组（$n=4$ 或 $n=5$），都不是严格双曲的。实际上，严格双曲性的失效常是对称性的副产品。例如，方程组\eqref{align:ch07-各向同性不可压缩}与\eqref{align:ch07-弹性弦}不是严格双曲的。

在规模 $n=2$ 的方程组中，严格双曲性通常在孤立的\emph{脐点（umbilic points）}处失效，在这些点处 $\mathrm{D}F$ 退化为单位矩阵的倍数。即使仅存在一个脐点，也足以在解的行为中造成混乱。这将在后续章节中通过下述简单方程组来展示：
\begin{align}
\begin{cases}
\partial_t u+\partial_x[(u^2+v^2)u]=0\\
\partial_t v+\partial_x[(u^2+v^2)v]=0,
\end{cases}
\label{align:ch07-脐点模型}
\end{align}
它是\eqref{align:ch07-各向同性不可压缩}与\eqref{align:ch07-弹性弦}的一个漫画式模型。\eqref{align:ch07-脐点模型}的特征速度为
\begin{align}
\lambda_1=u^2+v^2,\qquad \lambda_2=3(u^2+v^2),
\label{align:ch07-脐点模型特征速度}
\end{align}
相应的特征向量为
\begin{align}
R_1=\begin{pmatrix}v\\-u\end{pmatrix},\qquad R_2=\begin{pmatrix}u\\v\end{pmatrix},
\label{align:ch07-脐点模型特征向量}
\end{align}
因而该方程组是严格双曲的，除原点 $(0,0)$ 之外，而原点是一个脐点。

我们以一个有用恒等式的推导来结束本节。对等式 $\mathrm{D}F R_j=\lambda_j R_j$ 的两边作用 $\mathrm{D}$，再从左乘 $R_k^\top$；我们还对 $\mathrm{D}F R_k=\lambda_k R_k$ 作用 $\mathrm{D}$，再从左乘 $R_j^\top$。合并所得两方程，我们推出
\begin{align}
(\mathrm{D}\lambda_j R_k)R_j-(\mathrm{D}\lambda_k R_j)R_k=\mathrm{D}F[R_j,R_k]-\lambda_j\mathrm{D}R_j R_k+\lambda_k\mathrm{D}R_k R_j,\qquad j,k=1,\cdots,n,
\label{align:ch07-Lie括号恒等式}
\end{align}
其中 $[R_j,R_k]$ 表示 Lie 括号：
\begin{align}
[R_j,R_k]=\mathrm{D}R_j R_k-\mathrm{D}R_k R_j.
\label{align:ch07-Lie括号定义}
\end{align}
特别地，在严格双曲性失效的点 $U\in\mathscr{O}$（譬如 $\lambda_j(U)=\lambda_k(U)$），\eqref{align:ch07-Lie括号恒等式}给出
\begin{align}
(\mathrm{D}\lambda_j R_k)R_j-(\mathrm{D}\lambda_k R_j)R_k=(\mathrm{D}F-\lambda_j I)[R_j,R_k].
\label{align:ch07-退化点恒等式}
\end{align}
从左用 $L_j(U)$ 与 $L_k(U)$ 乘\eqref{align:ch07-退化点恒等式}，由\eqref{align:ch07-左右向量规范化}我们得出
\begin{align}
\mathrm{D}\lambda_j(U)R_k(U)=\mathrm{D}\lambda_k(U)R_j(U)=0.
\label{align:ch07-退化点Dlambda}
\end{align}

\end{document}
